% Maßstab, Ähnlichkeit und Strahlensätze – bank/strahlensaetze/ (zusammenbau v0.4, Heft)
\einheitenkopf[t8]{Maßstab, Ähnlichkeit und Strahlensätze}
\setcounter{aufgabe}{57}

% Anlauf (ohne Prüfkennung)
\begin{aufgabe}{Maßstab umrechnen – Anlauf}
\begin{teile}
\teil Maßstab $1 : 20$ – ist die Zeichnung kleiner oder größer als die Wirklichkeit? \kreuz{kleiner} \\ \kreuz{größer}
\teil Maßstab $1 : 20$ – wie viel ist $1\,\text{cm}$ auf dem Plan in Wirklichkeit? \leerfeld[cm]
\teil Ein Tisch ist $150\,\text{cm}$ lang. Maßstab $1 : 25$ – wie lang ist er auf dem Plan? \leerfeld[cm]
\end{teile}
\end{aufgabe}

\begin{aufgabe}{Maßstab umrechnen – Prüfungsaufgaben}
\begin{teile}
\teil Ein Leuchtturm ist $28{,}0\,\text{m}$ hoch und hat unten einen Durchmesser von $5{,}6\,\text{m}$. Für eine Ausstellung wird ein Modell im Maßstab $1 : 40$ gebaut. Berechne die Höhe und den Durchmesser des Modells in Zentimetern. \hfill (P10 2018 OS) \\ Höhe \leerfeld[cm] , Durchmesser \leerfeld[cm]
\teil Ein Wasserturm ist $36{,}0\,\text{m}$ hoch und hat einen Durchmesser von $9{,}6\,\text{m}$. Ein Modell wird im Maßstab $1 : 80$ gebaut. Berechne die Höhe und den Durchmesser des Modells in Zentimetern. \hfill (P10 2018 OS) \\ Höhe \leerfeld[cm] , Durchmesser \leerfeld[cm]
\teil Das Modell einer Kirche im Maßstab $1 : 25$ ist $0{,}92\,\text{m}$ hoch und $0{,}64\,\text{m}$ breit. Berechne die Höhe und die Breite der echten Kirche. \hfill (P10 2015 OS) \\ Höhe \leerfeld[m] , Breite \leerfeld[m]
\teil Das Modell eines Segelschiffs im Maßstab $1 : 30$ ist $0{,}45\,\text{m}$ lang und $0{,}62\,\text{m}$ hoch. Berechne die Länge und die Höhe des echten Schiffs. \hfill (P10 2015 OS) \\ Länge \leerfeld[m] , Höhe \leerfeld[m]
\end{teile}
\end{aufgabe}

% Anlauf (ohne Prüfkennung)
\begin{aufgabe}{Maßstabsgerecht zeichnen – Anlauf}
\begin{teile}
\teil Ein Beet ist $4\,\text{m}$ lang und $2\,\text{m}$ breit. Das Zeichenfeld ist $15\,\text{cm}$ breit und $9\,\text{cm}$ hoch. Welcher Maßstab passt? \kreuz{$1 : 20$} \\ \kreuz{$1 : 40$} \\ \kreuz{$1 : 1\,000$}
\teil Im Raster steht ein Kästchen für $2\,\text{m}$. Wie viele Kästchen lang ist eine $14\,\text{m}$ lange Mauer? Trage die Strecke im Raster ab.

\begin{ksys}[xmin=0,xmax=10,ymin=0,ymax=3]\end{ksys}
\leerfeld[Kästchen]
\teil Das Rechteck $ABCD$ ist $4$ Kästchen breit und $3$ Kästchen hoch. Zeichne es doppelt so groß.

\begin{ksys}[xmin=0,xmax=14,ymin=0,ymax=10]\punkt{1}{1}{A}\punkt{5}{1}{B}\punkt{5}{4}{C}\punkt{1}{4}{D}\end{ksys}
\end{teile}
\end{aufgabe}

\begin{aufgabe}{Maßstabsgerecht zeichnen – Prüfungsaufgaben}
\begin{teile}
\steil Eine zylinderförmige Futtertonne ist $90\,\text{cm}$ hoch und hat den Radius $33\,\text{cm}$. Sie soll mit einer rechteckigen Holzplatte abgedeckt werden, die $60\,\text{cm}$ breit und $85\,\text{cm}$ lang ist. Zeichne die Draufsicht von Tonne und Platte in einem passenden Maßstab in das Karofeld, gib den Maßstab an und beschrifte. Entscheide, ob die Platte die Tonne vollständig abdeckt. \hfill (P10 2021 OS) \\

\begin{ksys}[xmin=0,xmax=13,ymin=0,ymax=14]\end{ksys}
Maßstab: \leerfeld
\steil Ein runder Gartenpool hat den Radius $1{,}5\,\text{m}$. Er soll mit einer rechteckigen Plane abgedeckt werden, die $2{,}8\,\text{m}$ breit und $4\,\text{m}$ lang ist. Zeichne die Draufsicht von Pool und Plane in einem passenden Maßstab in das Karofeld, gib den Maßstab an und beschrifte. Entscheide, ob die Plane den Pool vollständig abdeckt. \hfill (P10 2021 OS) \\

\begin{ksys}[xmin=0,xmax=13,ymin=0,ymax=14]\end{ksys}
Maßstab: \leerfeld
\end{teile}
\end{aufgabe}
